% 固定假设的集中与数据依赖选择之间的区别。
\begin{tikzpicture}[x=1mm,y=1mm,font=\footnotesize]
  \draw[book arrow] (0,0) -- (94,0) node[right] {假设};
  \draw[book arrow] (0,0) -- (0,48) node[above] {分类错误率};
  \draw[BookRule,thin] (0,10) -- (92,10);
  \draw[BookRule,thin] (0,25) -- (92,25);
  \draw[BookRule,thin] (0,40) -- (92,40);

  \foreach \x/\pop/\emp in {10/31/29,25/22/19,40/27/23,55/18/12,70/24/20,85/20/17}{
    \draw[FigureOutput,very thick] (\x-3,\pop) -- (\x+3,\pop);
    \fill[FigureInput] (\x,\emp) circle (1.3);
    \draw[BookMuted,thin] (\x,\emp) -- (\x,\pop);
  }
  \node[font=\scriptsize,text=FigureOutput,anchor=west] at (5,45) {横线：期望风险$R_{\mathcal P}(h)$};
  \node[font=\scriptsize,text=FigureInput,anchor=west] at (47,45) {圆点：经验风险$\hat R_S(h)$};
  \draw[FigureLoss,very thick,rounded corners] (49,7) rectangle (61,34);
  \node[font=\kaishu\scriptsize,text=FigureLoss,align=center] at (55,4)
    {ERM选择了经验风险最低的假设};
  \draw[book arrow,FigureLoss] (55,4.5) -- (55,10.5);
  \node[font=\kaishu\scriptsize,text=BookMuted,align=center] at (46,-9)
    {一致收敛要求所有竖线同时很短，\\这样数据选出的假设也仍受控制。};
\end{tikzpicture}
